\section{Related Work}
\label{sec:related_work}

To contextualize our approach, we reviewed studies focused on gait cycle analysis and classification, particularly in post-stroke or pathological gait conditions.
Yuwono et al. \cite{ref10} proposed a method for classifying healthy and pathological gait using Shimmer IMUs placed on the legs. They extracted spectral features 
from the Short-Time Fourier Transform (STFT) of gyroscope signals and achieved a classification accuracy of 97\% using a Self-Organizing Map (SOM) combined with 
k-means clustering. Building upon this, Zhou et al. \cite{ref11} developed an adaptive gait event detection algorithm using a smart insole equipped with pressure 
sensors and IMUs.

Further advancements were introduced by Tunca et al. \cite{ref12}, who proposed a smartphone- and smartwatch-based system that utilized raw IMU data without preprocessing. 
Their deep learning model, DeepSense, combined Convolutional Neural Networks (CNNs) and Gated Recurrent Units (GRUs) to form a general-purpose framework for human activity 
recognition, achieving approximately 94\% classification accuracy. Mellone et al.\ \cite{ref13} validated the use of smartphone-based inertial sensors to capture gait parameters 
during the Timed Up and Go test. Their results support the feasibility of low-cost remote gait assessments, providing a foundation for non-laboratory clinical use. In a broader 
review, Bai et al.\ \cite{ref14} highlighted the variety of wearable sensor technologies and feature extraction techniques used for gait recognition, emphasizing their applications 
in health monitoring.

Rojek et al.\ \cite{ref15} focused on AI-based methods in post-stroke gait analysis, exploring deep learning and hybrid systems for classifying and interpreting gait 
cycles. Zhou et al.\ \cite{ref16} further improved gait recognition using a single waist-worn IMU by combining spatial and temporal feature extraction through a CNN 
and Bi-LSTM architecture, achieving 98.63\% accuracy. Panah et al.\ \cite{ref17} conducted a comprehensive review of wearable sensors in rehabilitation, outlining their 
roles in mobility assessment and feedback-based training. Muro-de-la-Herran et al.\ \cite{ref18} reviewed both wearable and non-wearable gait analysis systems, comparing 
their clinical applications, limitations, and use cases.

Ferrante et al.\ \cite{ref20} examined the effect of biofeedback cycling rehabilitation on gait recovery in 153 chronic stroke patients. Using 3D motion capture and 
gait pattern classification, they reported improvements in 51\% of participants. Finally, Kaczmarczyk et al.\ \cite{ref21} employed an artificial neural network (ANN) 
to classify post-stroke gait patterns from joint angle trajectories captured via the APAS motion system. Their model achieved classification accuracies of up to 100\% 
for knee joint data and approximately 92.5\% overall.

These studies collectively illustrate the wide range of sensing technologies and learning models applied to gait cycle analysis. However, despite these advances, 
there remains a notable gap in lightweight, interpretable methods that detect clinically meaningful gait asymmetries—particularly in post-stroke populations with 
limited, real-world sensor data.
